\section{Logical and computational foundations}
\label{sec:foundations}

\paragraph{Intrinsic multi-agent games.}
Abramsky's $L_A$ treats player-owned connectives, quantifiers, permutations,
and information constraints as constructors of interaction
\cite{Abramsky2007}. The public scripts here select a finite dependent
fragment of that viewpoint. Packaging the terminal objectives as the
per-agent valuation of one interpreted terminal gives the same outcome
evaluation for the corresponding first-order prefix. The larger concurrent
semantics is a source of constructions for extending this fragment.

Olde Loohuis and Venema study propositional and modal player-indexed logics
with arbitrary winner subsets at atoms, deriving equational and representation
results \cite{OldeLoohuisVenema2010}. Their winning-strategy equivalence is a
useful observation of games. Our execution interface retains the trace and
endpoint returns, since later typed continuations can depend on them.

\paragraph{Multirole proof theory and coherence.}
Multirole logic interprets formulas at sets of roles and proves multiparty
cut elimination \cite{XiWu2023}. It provides a natural foundation for extending
the singleton endpoint annotation $G@i$ to $G@R$. Multiparty session
coherence gives another account of how several endpoints jointly implement
one interaction \cite{CarboneEtAl2016}. These constructions guide the open
interface problem in \cref{sec:development}: its typing evidence must describe
a compatible collection, not just pairs declared dual by convention.

The public-core interaction theorem is a concrete witness of this
compatibility for one scripted connection. It should not be read as a
transfer of every cut or identity rule from either established system. The
terminal families and the interfaces required by each rule must be matched.

\paragraph{Dependent communication.}
Dependent session types interpret input and output through quantification
over values appearing in continuation types
\cite{ToninhoCairesPfenning2011}. Wu and Xi's multiparty dependent system
contains a role-indexed quantifier with endpoint-specific interpretation
\cite{WuXi2018}. Equation (4) uses the owner-as-witness convention to expose
that mechanism in a small data calculus. Its operational content is the
shared substitution in (14).

Dependent game trees and the computation of optimal strategies are developed
by Escard\'o and Oliva \cite{EscardoOliva2023}. Compositional game theory
similarly studies games as open interacting components
\cite{GhaniEtAl2018}. These results provide complementary semantic tools:
dependent choice describes the available actions, while an equilibrium or
optimality condition describes a property of selected implementations.

\paragraph{Availability and modal structure.}
A private value and a public value can have the same underlying data type
while admitting different uses. This is a judgmental issue: it concerns the
context from which a term may be built. Dependency calculi and multimodal
dependent type theory supply principled accounts of such distinctions
\cite{ChoudhuryEadesWeirich2022,GratzerEtAl2021}. Structural information flow
also connects polarity, dependencies, and controlled disclosure
\cite{GouniEtAl2025}. An extension of (20)--(22) should explain its context
rules and their interaction with dependent conversion using that vocabulary.

\paragraph{Control and dependent data.}
The $\lambda\mu\tilde\mu$ calculus exposes the program/context and
call-by-name/call-by-value symmetries of classical computation
\cite{CurienHerbelin2000}. A multi-role control calculus could give another
operational account of interaction, but control adds computational structure
beyond the binder in (2).

Herbelin shows that particular unrestricted combinations of computational
control, strong dependent sums, equality, and conversion collapse the
quantification domain \cite{Herbelin2005}. This is a constraint on the
combination of rules. Miquey develops a dependent classical sequent calculus
using restricted dependencies and delimited continuations
\cite{Miquey2017}. The public core uses stable data and ordinary dependent
elimination, so no control principle is required for its interaction law.

\paragraph{Position of the construction.}
The common thread is ownership-sensitive binding. In the public core, it has
both a global script and a sum/product endpoint interpretation. In the
information semantics, the same chosen value persists while its audience
changes. In a larger proof calculus, cut must compose these bindings without
changing their types, owners, or permitted dependencies. These are connected
requirements on one interaction object, and each has an explicit mathematical
interface in the development above.
